% Conclusion for the Introspective Verification paper.
% Style: plain, direct; no em-dashes; minimal colons/quotes; no anthropomorphic verbs.

\section{Conclusion}
\label{sec:conclusion}

We asked whether the judgment that a generative process reward model produces by decoding a critique is already available in the verifier's hidden states, without generation. It is. Step correctness is decodable by a probe on an intermediate layer of the backbone, near seventy percent of the depth, more strongly than from the final layer, because the final state is shaped to emit the surface verdict while an earlier state reflects whether the step follows from its premises. We turned this into an introspective verdict head, a residual readout over an intermediate and a final layer, trained with LoRA on step labels and evaluated in a single forward pass on a single V100. On ProcessBench at 1.5B it reaches 58.1 mean F1, above most trained process reward models including ones at 7B and 14B scale and above a same size generative model, from a read that is about an order of magnitude faster in wall clock than a generative pass. The two signals are complementary, and a fixed combination with a generative verifier improves it at both scales.

The picture we report is deliberately honest about where the cheap read stops. At matched training the residual form does not raise the mean over the strongest single readout; it redistributes accuracy toward the proof heavy subsets while the intermediate view alone is the weaker of the two. The per domain calibration that helps the 7B head hurts the 1.5B head, so we present it as a scale specific trick rather than a component of the method. The head is over confident and needs a tuned threshold, its localization is coarser than an explicit re-derivation on the hardest subsets, and it cannot catch errors of omission, which are invisible to any step local read. These are limits of reading a single forward pass, and they mark the regime where generation earns its cost.

Two directions follow. The first is downstream use, since a verifier is ultimately measured by the answers it selects, and an introspective reward that costs one read is attractive for Best of N selection and for search. The second is to close the localization gap on proof heavy problems, either by reading more than one intermediate layer or by spending a short generation only on the steps where the head is uncertain, which keeps the average cost near one read while recovering the cases where generation currently wins. More broadly, the result suggests that verification, like other judgments a model can make about its own computation, may not need to be spoken to be read.
